\begin{helper}
Let \(g:\mathbb R^3\to\mathbb R\) be continuous, written \(g(x,y,z)\), and
assume \(g(x,y,z)\ge0\) whenever \(0\le x\le y\le z\le1\).  Then the following
regularity facts hold.

First, for every \(0\le y\le z\le1\), the function
\(x\mapsto g(x,y,z)\) is interval-integrable on \([0,y]\) and is nonnegative
on \((0,y]\).

Second, for every \(0\le z\le1\), the function
\[
F_z(y)=\int_0^y g(x,y,z)\,dx
\]
is continuous on \([0,z]\), interval-integrable on \([0,z]\), and
nonnegative on \((0,z]\).

Third, the function
\[
G(z)=\int_0^z\int_0^y g(x,y,z)\,dx\,dy
\]
is continuous on \([0,1]\), interval-integrable on \([0,1]\), and
nonnegative on \((0,1]\).
\end{helper}
\begin{proof}
Fix \(0\le y\le z\le1\).  The map \(x\mapsto g(x,y,z)\) is continuous because
it is the composition of the continuous map \(x\mapsto(x,y,z)\) with the
continuous function \(g\).  A continuous real-valued function is
interval-integrable on every compact interval, so it is interval-integrable on
\([0,y]\).  If \(x\in(0,y]\), then \(0\le x\le y\le z\le1\), and the simplex
nonnegativity hypothesis gives \(g(x,y,z)\ge0\).

Now fix \(0\le z\le1\).  Consider
\[
F_z(y)=\int_0^y g(x,y,z)\,dx .
\]
The two-variable map \((x,y)\mapsto g(x,y,z)\) is continuous.  The theorem on
continuity of parameter-dependent interval integrals with continuous
integrand and continuous moving endpoint therefore gives continuity of
\(F_z\) as a function of \(y\); restricting to \([0,z]\) gives continuity on
that compact interval.  Hence \(F_z\) is interval-integrable on \([0,z]\).  If
\(y\in(0,z]\), then the first paragraph applies to every \(x\in(0,y]\), so
the integrand in \(\int_0^y g(x,y,z)\,dx\) is nonnegative almost everywhere on
\((0,y]\).  Monotonicity of the real integral over an interval with
nonnegative orientation gives \(F_z(y)\ge0\).

Finally define
\[
G(z)=\int_0^z F_z(y)\,dy
=\int_0^z\int_0^y g(x,y,z)\,dx\,dy .
\]
The preceding parameter-integral continuity theorem, applied to the
continuous three-variable integrand \(g\) and the continuous endpoint
\((y,z)\mapsto y\), gives continuity of \((y,z)\mapsto F_z(y)\) on the
triangle \(0\le y\le z\le1\).  Applying the same theorem once more to the
outer integral with endpoint \(z\) gives continuity of \(G\) on \([0,1]\).
Thus \(G\) is interval-integrable on \([0,1]\).  If \(z\in(0,1]\), the second
paragraph gives \(F_z(y)\ge0\) for every \(y\in(0,z]\), so real-integral
monotonicity on \([0,z]\) gives \(G(z)\ge0\).
\end{proof}
